
\subsubsection{3.1 Overview and Hypothesis Map}

The diagnostic approach rests on a single observation: if backbone-level spurious encoding and head-level decision making are the two sites of robustification, then a measurement sensitive to one and not the other is needed. A linear probe trained on \textit{frozen} backbone features to predict the \textit{spurious attribute} is exactly this measurement --- it captures what the backbone encodes while remaining independent of head weights.

Five pre-registered hypotheses together constitute a three-gate verification study:

\begin{table}[htbp]
\centering
\resizebox{\columnwidth}{!}{%
\begin{tabular}{lll}
\toprule
Gate ID & Substantive Claim & Gate Type \\
\midrule
H-P0 & DFR backbone = ERM backbone at matching seeds & MUST\_WORK \\
H-M1 & GroupDRO minority group upweighting creates group-balanced gradient signal & MUST\_WORK \\
H-M2 & Group-balanced gradient propagates to all backbone layers & SHOULD\_WORK \\
H-M3 & GroupDRO reduces background linear decodability in layer4 features & MUST\_WORK (primary) \\
H-P2 & Background probe accuracy correlates negatively with WGA & SHOULD\_WORK (exploratory) \\
\bottomrule
\end{tabular}
}
\end{table}

Gates are declared CONFIRMED (p < 0.05, d > 0), SUGGESTIVE (p < 0.10, d > 0.5), or REJECTED based on pre-registered thresholds [see \texttt{03\textbackslash \_refinement.yaml}].

\subsubsection{3.2 Checkpoints and Dataset}

\textbf{Checkpoints.} The publicly available izmailovpavel/spurious\_feature\_learning checkpoints \citep{izmailov2022feature} provide 12 ResNet-50 models: 3 seeds $\times$ 4 training methods (ERM, SAM, GroupDRO, DFR). Matched seeds across all 4 methods enable paired statistical tests that control within-seed variance.

\textbf{Dataset.} Waterbirds WILDS [Sagawa et al., 2019; Koh et al., 2021] provides 4,795 training and 5,794 test images. The \texttt{group\textbackslash \_array \% 2} annotation provides binary background labels (0 = land background, 1 = water background). Minority groups (landbird on water, waterbird on land) constitute 5.01\% of the training set (240 of 4,795 images).

\subsubsection{3.3 Background Linear Probe Protocol}

\textbf{Feature extraction.} For each checkpoint, layer4 features are extracted using a forward hook with \texttt{torch.no\textbackslash \_grad()}, followed by \texttt{AdaptiveAvgPool2d(output\textbackslash \_size=(1,1))} to produce D=2048 feature vectors, evaluated on the full test set (N=5,794).

\textbf{Probe classifier.} \texttt{sklearn.linear\textbackslash \_model.LogisticRegression(solver='lbfgs', C=1e9, max\textbackslash \_iter=1000, random\textbackslash \_state=42)} is trained to predict background attribute (land vs. water) from layer4 features. C=1e9 (effectively no regularization) follows the established convention [Kirichenko et al., 2022; Murotkar et al., 2024] and avoids optimizer geometry confounds.

\subsubsection{3.4 DFR Backbone Identity Verification (H-P0)}

Mean cosine similarity of layer4 feature vectors is computed for DFR vs. ERM checkpoint pairs at matching seeds (50 test images, 3 seed pairs). \textbf{Pre-registered gate:} mean cosine similarity $\geq$ 0.9999 for all 3 seed pairs.

\subsubsection{3.5 Gradient Propagation Verification (H-M2)}

L2 norms of weight differences (GroupDRO $-$ ERM) per ResNet-50 block are computed, and the backbone-to-head ratio is reported. $DFR-ERM$ serves as a negative control. \textbf{Pre-registered gate:} backbone/head ratio > 1 for all 3 seeds.

The H-M2 linear probe (fit on a smaller probe training set than H-M3) yields p = 0.0526 (SUGGESTIVE), while the primary H-M3 probe using the full N=5,794 test set yields p = 0.0039. The difference in statistical confidence likely reflects the larger and more stable probe accuracy estimates from N=5,794 samples.

\subsubsection{3.6 Primary Probe Accuracy Test (H-M3)}

A one-sided paired t-test (GroupDRO < ERM direction, n = 3 seed pairs) is applied to background probe accuracy computed on the full test set (N=5,794). \textbf{Pre-registered thresholds:} CONFIRMED if p < 0.05 and d > 0; SUGGESTIVE if 0.05 $\leq$ p < 0.10 and d > 0.5.

\subsubsection{3.7 WGA Correlation Analysis (H-P2, Exploratory)}

Pearson r between background probe accuracy and WGA [Izmailov et al., 2022, Table 1] is computed across 9 checkpoints (ERM $\times$ 3, SAM $\times$ 3, GroupDRO $\times$ 3). DFR is excluded because its backbone is identical to ERM (H-P0), making its probe accuracy indistinguishable from ERM's. A 95\% bootstrap confidence interval (percentile method, n\_resamples=1,000) is reported. BCa bootstrap was not used because it produces degenerate (NaN) bounds at n=9 due to insufficient distinct samples; 2 of 1,000 bootstrap resamples yielded NaN and were excluded, leaving 998 clean samples for CI computation. \textbf{Pre-registered gate (SHOULD\_WORK):} r < $-0.5$.

WGA values from Izmailov et al. [2022] are method-level constants (ERM=0.72, SAM=0.74, GroupDRO=0.88) with no per-seed variation reported. This creates within-method WGA collinearity that inflates bootstrap CI width, limiting the analysis to SUGGESTIVE at n=9.

\subsubsection{3.8 Mechanistic Chain Structure}

The five hypotheses form a verified mechanistic chain with DFR establishing the structural alternative:

\begin{verbatim}
GroupDRO minority upweighting (H-M1, MUST_WORK)
    → Group-balanced gradient propagates to backbone (H-M2, SHOULD_WORK)
        → Reduced background linear decodability in layer4 (H-M3, MUST_WORK)
            → Correlates with improved WGA (H-P2, SHOULD_WORK, exploratory)

DFR: head recalibration on unchanged backbone features (H-P0, backbone cosine sim = 1.000000)
    → WGA = 0.91 without any backbone modification
\end{verbatim}

